\documentclass{amsart}
% For dealing with references we use the comment environment
\usepackage{verbatim}
\newenvironment{reference}{\comment}{\endcomment}
%\newenvironment{reference}{}{}
\newenvironment{slogan}{\comment}{\endcomment}
\newenvironment{history}{\comment}{\endcomment}

% For commutative diagrams we use Xy-pic
\usepackage[all]{xy}

% We use 2cell for 2-commutative diagrams.
\xyoption{2cell}
\UseAllTwocells

% We use multicol for the list of chapters between chapters
\usepackage{multicol}

% This is generally recommended for better output
\usepackage{lmodern}
\usepackage[T1]{fontenc}

% For cross-file-references

% Package for hypertext links:
\usepackage{hyperref}

% For any local file, say "hello.tex" you want to link to please

% Theorem environments.
%
\theoremstyle{plain}
\newtheorem{theorem}[subsection]{Theorem}
\newtheorem{proposition}[subsection]{Proposition}
\newtheorem{lemma}[subsection]{Lemma}

\theoremstyle{definition}
\newtheorem{definition}[subsection]{Definition}
\newtheorem{example}[subsection]{Example}
\newtheorem{exercise}[subsection]{Exercise}
\newtheorem{situation}[subsection]{Situation}

\theoremstyle{remark}
\newtheorem{remark}[subsection]{Remark}
\newtheorem{remarks}[subsection]{Remarks}

\numberwithin{equation}{subsection}

% Macros
%
\def\lim{\mathop{\mathrm{lim}}\nolimits}
\def\colim{\mathop{\mathrm{colim}}\nolimits}
\def\Spec{\mathop{\mathrm{Spec}}}
\def\Hom{\mathop{\mathrm{Hom}}\nolimits}
\def\Ext{\mathop{\mathrm{Ext}}\nolimits}
\def\SheafHom{\mathop{\mathcal{H}\!\mathit{om}}\nolimits}
\def\SheafExt{\mathop{\mathcal{E}\!\mathit{xt}}\nolimits}
\def\Sch{\mathit{Sch}}
\def\Mor{\mathop{\mathrm{Mor}}\nolimits}
\def\Ob{\mathop{\mathrm{Ob}}\nolimits}
\def\Sh{\mathop{\mathit{Sh}}\nolimits}
\def\NL{\mathop{N\!L}\nolimits}
\def\CH{\mathop{\mathrm{CH}}\nolimits}
\def\proetale{{pro\text{-}\acute{e}tale}}
\def\etale{{\acute{e}tale}}
\def\QCoh{\mathit{QCoh}}
\def\Ker{\mathop{\mathrm{Ker}}}
\def\Im{\mathop{\mathrm{Im}}}
\def\Coker{\mathop{\mathrm{Coker}}}
\def\Coim{\mathop{\mathrm{Coim}}}

% Boxtimes
%
\DeclareMathSymbol{\boxtimes}{\mathbin}{AMSa}{"02}

%
% Macros for moduli stacks/spaces
%
\def\QCohstack{\mathcal{QC}\!\mathit{oh}}
\def\Cohstack{\mathcal{C}\!\mathit{oh}}
\def\Spacesstack{\mathcal{S}\!\mathit{paces}}
\def\Quotfunctor{\mathrm{Quot}}
\def\Hilbfunctor{\mathrm{Hilb}}
\def\Curvesstack{\mathcal{C}\!\mathit{urves}}
\def\Polarizedstack{\mathcal{P}\!\mathit{olarized}}
\def\Complexesstack{\mathcal{C}\!\mathit{omplexes}}
% \Pic is the operator that assigns to X its picard group, usage \Pic(X)
% \Picardstack_{X/B} denotes the Picard stack of X over B
% \Picardfunctor_{X/B} denotes the Picard functor of X over B
\def\Pic{\mathop{\mathrm{Pic}}\nolimits}
\def\Picardstack{\mathcal{P}\!\mathit{ic}}
\def\Picardfunctor{\mathrm{Pic}}
\def\Deformationcategory{\mathcal{D}\!\mathit{ef}}

\setlength{\emergencystretch}{3em}
\begin{document}
\title[Stacks Project, Tag 05EK]{409 --- Exact recovery complex for modules under flat neighbourhood gluing}
\maketitle
\noindent Copyright (C) 2005--2025 Johan de Jong. Stacks Project authors. Source: \href{https://stacks.math.columbia.edu/tag/05EK}{Tag 05EK}.

\noindent Editorial difficulty note (not author text). Difficulty H2: The author proves injectivity using ideal-power torsion and the unchanged infinitesimal neighbourhood. Its computational converse tracks a compatible tuple through a Koszul-homology obstruction, uses flatness to identify that obstruction with its base change, and recovers the tuple from one module element..

\noindent Editorial provenance note (not author text). History: excerpt from revision \texttt{a0444\allowbreak 6e57e\allowbreak c1fbc\allowbreak 252a8\allowbreak 71afc\allowbreak ec775\allowbreak 2fb28\allowbreak 07b14}, more-algebra.tex, lines 24578--24638. Reference presentation was adapted; original-author context and core support blocks were retained or added. The mathematical wording is unchanged. Licensed under GNU FDL 1.2 or later; no invariant sections or cover texts. See ../licenses/COPYING. Context blocks reproduce author definitions and supporting results; they are not additional counted problems.

Let $R \to S$ be a ring map.
Let $f_1, \ldots, f_t \in R$ and $I = (f_1, \ldots, f_t)$.
Then for any $R$-module $M$ we can define a complex
\begin{equation}
\label{equation-glueing-complex}
0 \to M \xrightarrow{\alpha}
M \otimes_R S \times \prod M_{f_i} \xrightarrow{\beta}
\prod (M \otimes_R S)_{f_i}
\times
\prod M_{f_if_j}
\end{equation}
where $\alpha(m) = (m \otimes 1, m/1, \ldots, m/1)$ and
$$
\beta(m', m_1, \ldots, m_t) =
((m'/1 - m_1 \otimes 1, \ldots, m'/1 - m_t \otimes 1),
(m_1 - m_2, \ldots, m_{t - 1} - m_t).
$$
We would like to know when this complex is exact.


\begin{lemma}
\label{lemma-recover-module-from-glueing-data}
Assume $\varphi : R \to S$ is a flat ring map and
$I = (f_1, \ldots, f_t) \subset R$ is an ideal such that
$R/I \to S/IS$ is an isomorphism.
Let $M$ be an $R$-module. Then the
complex (\ref{equation-glueing-complex})
is exact.
\end{lemma}

\begin{proof}
First proof. Denote $\check{\mathcal{C}}_R \to \check{\mathcal{C}}_S$
the quasi-isomorphism of extended alternating {\v C}ech complexes of
Lemma \href{https://stacks.math.columbia.edu/tag/091M}{lemma map identifies koszul and cech complexes (Tag 091M)}.
Since these complexes are bounded with flat terms, we see that
$M \otimes_R \check{\mathcal{C}}_R \to M \otimes_R \check{\mathcal{C}}_S$
is a quasi-isomorphism too (Lemmas
\href{https://stacks.math.columbia.edu/tag/064K}{lemma derived tor quasi isomorphism (Tag 064K)} and
\href{https://stacks.math.columbia.edu/tag/064L}{lemma derived tor quasi isomorphism other side (Tag 064L)}). Now the complex
(\ref{equation-glueing-complex}) is a truncation of the cone
of the map
$M \otimes_R \check{\mathcal{C}}_R \to M \otimes_R \check{\mathcal{C}}_S$
and we win.

\medskip\noindent
Second computational proof.
Let $m \in M$. If $\alpha(m) = 0$, then $m \in M[I^\infty]$, see
Lemma \href{https://stacks.math.columbia.edu/tag/05EA}{lemma torsion free (Tag 05EA)}. Pick $n$ such that $I^n m = 0$
and consider the map $\varphi : R/I^n \to M$.
If $m \otimes 1 = 0$, then $\varphi \otimes 1_S = 0$, hence
$\varphi = 0$ (see
Lemma \href{https://stacks.math.columbia.edu/tag/05EC}{lemma neighbourhood equivalence (Tag 05EC)})
hence $m = 0$. In this way we see that $\alpha$ is injective.

\medskip\noindent
Let $(m', m'_1, \ldots, m'_t) \in \Ker(\beta)$.
Write $m'_i = m_i/f_i^n$ for some $n > 0$ and $m_i \in M$.
We may, after possibly enlarging $n$ assume that
$f_i^n m' = m_i \otimes 1$ in $M \otimes_R S$ and
$f_j^nm_i - f_i^nm_j = 0$ in $M$.
In particular we see that
$(m_1, \ldots, m_t)$ defines an element $\xi$ of
$H_1(M, (f_1^n, \ldots, f_t^n))$.
Since $H_1(M, (f_1^n, \ldots, f_t^n))$ is annihilated by $I^{tn + 1}$ (see
Lemma \href{https://stacks.math.columbia.edu/tag/05EF}{lemma koszul homology annihilated (Tag 05EF)})
and since $R \to S$ is flat we see that
$$
H_1(M, (f_1^n, \ldots, f_t^n)) =
H_1(M, (f_1^n, \ldots, f_t^n)) \otimes_R S =
H_1(M \otimes_R S, (f_1^n, \ldots, f_t^n))
$$
by
Lemma \href{https://stacks.math.columbia.edu/tag/05E9}{lemma neighbourhood isomorphism (Tag 05E9)}
The existence of $m'$ implies that $\xi$ maps to zero in the last group, i.e.,
the element $\xi$ is zero. Thus there exists an $m \in M$ such that
$m_i = f_i^n m$. Then $(m', m'_1, \ldots, m'_t) - \alpha(m)
= (m'', 0, \ldots, 0)$ for some $m'' \in (M \otimes_R S)[(IS)^\infty]$.
By
Lemma \href{https://stacks.math.columbia.edu/tag/05EC}{lemma neighbourhood equivalence (Tag 05EC)}
we conclude that $m'' \in M[I^\infty]$ and we win.
\end{proof}
\end{document}
